406
\hfill 6 Integration
if the integral on the right-hand side exists. This limit is called, following Cauchy,
the principal value of the integral, and, to distinguish the definitions (6.81) and
(6.82), we put the letters PV in front of the second to indicate that it is the principal
value.

In accordance with this convention we have
\textbf{Example 19}
$$
\text{PV} \int_{-1}^{1} \frac{dx}{x} = 0.
$$
We also adopt the following definition:
$$
\text{PV} \int_{-\infty}^{+\infty} f(x)dx := \lim_{R \to +\infty} \int_{-R}^{R} f(x)dx. \quad (6.83)
$$

\textbf{Example 20}
$$
\text{PV} \int_{-\infty}^{+\infty} x dx = 0.
$$

Finally, if there are several (finitely many) singularities of one kind or another on
the interval of integration, at interior points or endpoints, then the nonsingular points
of the interval are divided into a finite number of such intervals, each containing
only one singularity, and the integral is computed as the sum of the integrals over
the closed intervals of the partition.

It can be verified that the result of such a computation is not affected by the
arbitrariness in the choice of a partition.

\textbf{Example 21} The precise definition of the logarithmic integral can now be written
as
$$
\text{li}x = \begin{cases} \int_{0}^{x} \frac{dt}{\ln t}, & \text{if } 0 < x < 1, \\ \text{PV} \int_{0}^{x} \frac{dt}{\ln t}, & \text{if } 1 < x. \end{cases}
$$
In the last case the symbol PV refers to the only interior singularity on the interval
]0, $x$[ , which is located at 1. We remark that in the sense of the definition in formula
(6.81) this integral is not convergent.

\subsection{6.5.4 Problems and Exercises}
\textbf{1.} Show that the following functions have the stated properties.
a) $\text{Si}(x) = \int_{0}^{x} \frac{\sin t}{t} dt$ (the sine integral) is defined on all of $\mathbb{R}$, is an odd function,
and has a limit as $x \to +\infty$.
b) $\text{si}(x) = \int_{x}^{\infty} \frac{\sin t}{t} dt$ is defined on all of $\mathbb{R}$ and differs from $\text{Si}x$ only by a
constant;